% ---------------------------------------------------------------------------
\subsubsection{The pipeline at a glance}
% ---------------------------------------------------------------------------

Everything funnels into one DataFrame, \code{df_metrics}, the single source of
truth for every \code{groupby} and every plot. Solid (teal) arrows are the data
path; dashed (grey) arrows inject the PDDL ground truth needed to score plans
(Figure~\ref{fig:pipeline}).

\begin{figure}[h]
\centering
\resizebox{\linewidth}{!}{%
\begin{tikzpicture}[
  font=\ttfamily\scriptsize,
  diskbox/.style={draw=linecol,fill=white,rounded corners=3pt,inner sep=4pt,
                  align=left,minimum height=10mm,line width=0.8pt},
  inkbox/.style={draw=ink,fill=ink,text=white,rounded corners=3pt,inner sep=4pt,
                 align=left,minimum height=10mm,line width=0.8pt},
  widebox/.style={draw=accent,fill=accent!6,rounded corners=3pt,inner sep=5pt,
                  align=left,line width=1pt},
  data/.style={-{Stealth[length=5pt]},accent,line width=1pt},
  gt/.style={-{Stealth[length=5pt]},muted,line width=0.9pt,dashed},
  node distance=5mm and 7mm,
]
\node[diskbox] (raw) {\textbf{raw/}\\[-1pt]\textcolor{muted}{\tiny gen + reasoning}};
\node[diskbox,right=of raw] (parsed) {\textbf{parsed/}\\[-1pt]\textcolor{muted}{\tiny raw./reasoning.}};
\node[diskbox,right=of parsed] (scored) {\textbf{scored/}\\[-1pt]\textcolor{muted}{\tiny VAL: solved}};
\node[diskbox,right=of scored] (tasks) {\textbf{tasks/}\\[-1pt]\textcolor{muted}{\tiny domain+pfile}};

\node[inkbox,below=9mm of parsed.south,anchor=north,xshift=6mm] (load)
  {load\_artifact\_rows()\\[-1pt]\textcolor{linecol}{\tiny walk parsed, join raw+scored}};
\node[inkbox,below=9mm of tasks.south,anchor=north] (index)
  {index\_pddl\_context()\\[-1pt]\textcolor{linecol}{\tiny domain\_info / problem\_info}};

\node[widebox,below=10mm of load.south west,anchor=north west,text width=0.86\linewidth] (df)
  {compute\_row\_metrics() $\rightarrow$ \textbf{df\_metrics}\\[-1pt]
   \textcolor{accent}{\tiny hallucination $\cdot$ PDDLSimulator(exec/PAS) $\cdot$ CoTaL $\cdot$ \_iter1 $\cdot$ \_iwsr\_contrib}};

\node[inkbox,below=8mm of df.south west,anchor=north west] (agg)
  {build\_aggregate\_tables()\\[-1pt]\textcolor{linecol}{\tiny overall/domain/rank/PS+CI}};
\node[diskbox,right=8mm of agg] (plots) {plots $\times$32};
\node[diskbox,right=6mm of plots] (report) {report.json};

\draw[data] (raw) -- (parsed);
\draw[data] (parsed) -- (scored);
\draw[data] (scored.south) |- (load.east);
\draw[data] (parsed.south) -- (load.north);
\draw[data] (tasks.south) -- (index.north);
\draw[data] (load.south) -- ([xshift=-30mm]df.north);
\draw[gt]   (index.south) -- ([xshift=30mm]df.north);
\draw[data] (df.south -| agg) -- (agg.north);
\draw[data] (agg) -- (plots);
\draw[data] (plots) -- (report);
\end{tikzpicture}%
}
\caption{The analysis pipeline. One DataFrame, \texttt{df_metrics}, is the single
source of truth; PDDL ground truth (dashed) is injected to score plans.}
\label{fig:pipeline}
\end{figure}

% ---------------------------------------------------------------------------
\subsubsection{How a row is built: the three inputs to \texorpdfstring{\texttt{df\_metrics}}{df\_metrics}}
% ---------------------------------------------------------------------------

The canonical table is produced by \code{compute_row_metrics()}. It does
\emph{not} read the filesystem and it does \emph{not} compute ground truth ---
both happen before it. By the time the row loop runs, three things already
exist, each with a structurally different role.

\begin{table}[h]\small
\begin{tabularx}{\linewidth}{@{}c p{2.7cm} p{3.2cm} >{\raggedright\arraybackslash}X@{}}
\toprule
\# & \textbf{Input} & \textbf{Produced by} & \textbf{Role in the row loop}\\
\midrule
1 & \code{df_raw} & \code{load_artifact_rows()} &
The skeleton: one row per instance with the artifact-derived columns and the
keys used to look up ground truth. It is \code{.copy()}-ed, then \emph{widened}
with computed columns.\\[3pt]
2 & \code{domain_info} / \code{problem_info} & \code{index_pddl_context()} &
The ground truth: two dict lookups (\code{d_info}, \code{p_info}) consumed per
row to score the plan. \emph{Not} stored as columns --- they are the reference
the metric functions compare against.\\[3pt]
3 & \code{warnings_out} & threaded list via \code{add_warning()} &
The diagnostic side-channel: a by-reference list accumulating structured
records whenever an input is missing or malformed. Never part of
\code{df_metrics}; embedded separately in the report.\\
\bottomrule
\end{tabularx}
\end{table}

\begin{bmkeybox}{What \code{df\_metrics} literally is}
\code{df_metrics = df_raw.copy()} horizontally concatenated with three blocks of
computed columns hallucination, precondition/executability (from
\code{PDDLSimulator}), and CoTaL plus two derived columns \code{_iter1} and
\code{_iwsr_contrib}. So input~1 supplies the columns we keep; input~2 is
\emph{consumed} to manufacture the new columns and then discarded; input~3 is a
parallel log, not a column.
\end{bmkeybox}

\paragraph{\texorpdfstring{\texttt{df\_raw}}{df\_raw}-the artifact skeleton}

\code{load_artifact_rows()} walks the canonical path
\texttt{parsed/\allowbreak<run>/\allowbreak<model>/\allowbreak<protocol>/\allowbreak<domain>/\allowbreak<difficulty>/\allowbreak<instance>.json} and for each instance, joins the matching \code{scored/} data. 
It emits exactly \emph{one row per problem instance}, collapsing the retry loop by keeping the
last non-empty plan as the representative \code{_actions} while preserving the full attempt lists as private columns.

\begin{lstlisting}
# public: identity + outcome
Model, Domain, Problem, Difficulty, Protocol, Run_id
Valid (= scored["solved"]), Length, Iterations (= scored["iterations_used"]), Chain_of_Thought
# private:consumed downstream, dropped from the report
_actions          # representative plan (last non-empty attempt)
_cot_text         # reasoning trace text
_parsed_attempts  _scored_attempts  _raw_attempts   # full per-iteration lists
_file_path  parsed_file_path
\end{lstlisting}

The keys that matter for the next input are \code{Domain}, \code{Difficulty} and
\code{Problem} that are together the exact triple used to retrieve per-instance ground
truth. The \code{Chain_of_Thought} flag is resolved from three sources (protocol
YAML, presence of reasoning text, presence of reasoning actions) so CoT-bearing
runs are never mislabelled.

\begin{bmkeybox}{Why one row per instance, not per attempt}
Keeping the grain at the \emph{instance} level means every aggregate (success
rate, hallucination, PS) is automatically an average over comparable units. The
iteration structure is not lost it survives in the private
\code{_*_attempts} columns and the \code{Iterations} count but it is opt-in,
read only by metrics that need it (FASR, IWSR, the iteration profile, CoTaL).
\end{bmkeybox}

\paragraph{The PDDL parsing layer: \texorpdfstring{\texttt{d\_info}}{d\_info} and \texorpdfstring{\texttt{p\_info}}{p\_info}}

Plans can only be scored against the formal task they were generated for, so the
framework parses the PDDL itself. Two small, deliberately permissive regex
parsers do this, a syntax error emits a warning never aborts the run. They
answer two different questions and return two differently-shaped dicts.

\code{parse_domain_pddl(domain_path)} reads one \code{domain.pddl} (the schema
shared by every instance in a domain) and returns the legal vocabulary plus
per-action templates:

\begin{lstlisting}
parse_domain_pddl(...) -> {
  "action_names": set[str],   # every legal action head: the hallucination whitelist
  "predicates":   set[str],   # legal predicate names
  "functions":    set[str],   # numeric fluent names: what is numeric, not propositional
  "schemas": {                # one entry PER action, the simulator's instruction sheet
     action_name: {
        "params":   [str, ...],   # ordered ?param variable names, for grounding
        "prec_raw": str,          # raw :precondition s-expression text (unparsed)
        "eff_raw":  str,          # raw :effect s-expression text (unparsed)
     }, ...
  },
  "path": str | None,         # None when the file was missing (empty stub)
}
\end{lstlisting}

Each field has a single downstream consumer: \code{action_names} is the
whitelist for action hallucination; \code{functions} tells the simulator which
heads are numeric fluents; and \code{schemas} is what the simulator grounds, 
\code{params} maps the action's arguments onto the variables inside the raw
\code{prec_raw}/\code{eff_raw} strings so preconditions can be checked and
effects applied.

\code{parse_problem_pddl(problem_path)} reads one instance file (its own object
set and starting state) and returns:

\begin{lstlisting}
parse_problem_pddl(...) -> {
  "objects":      set[str],               # legal argument tokens: object-hallucination whitelist
  "init_atoms":   set[tuple[str, ...]],   # initial true propositions, e.g. ("at","truck1","depot")
  "init_numeric": dict[tuple, float],     # initial fluent values, e.g. ("fuel","truck1") -> 50.0
  "path":       str,
  "difficulty": str,                      # parent folder name
}
\end{lstlisting}

Object hallucination uses \code{objects}; the simulator seeds its world from
\code{init_atoms} and \code{init_numeric}. Action vs object-hallucination are
kept as separate signals on purpose, a model can know the legal action names
yet still invent objects which is a different and shallower grounding failure.

\code{index_pddl_context(tasks_dir)} calls those two parsers exactly once each
per file, walking \code{tasks/} a single time, and packs the results into two
lookups whose key shapes mirror the granularity of what they hold:

\begin{lstlisting}
domain_info  : dict[str, dict]                  # key = domain_name  (schema is shared)
problem_info : dict[tuple[str,str,str], dict]   # key = (domain, difficulty, instance)
\end{lstlisting}

This is the ``parse once, consume many'' design: every PDDL file is read once
regardless of how many rows or retries reference it. In the row loop the
retrieval is then two O(1) dict gets:

\begin{lstlisting}
d_info = domain_info.get(domain, {})                            # by domain
p_info = problem_info.get((domain, difficulty, instance), {})   # by full triple
pddl_ok = bool(d_info.get("action_names")) and bool(p_info.get("objects"))
\end{lstlisting}

\begin{bmcavbox}{A missing domain and a missing problem fail differently}
When \code{domain.pddl} is absent, \code{index_pddl_context} still stores a
populated-but-empty stub under the domain key (empty \code{action_names}, empty
\code{schemas}, \code{path=None}). When a problem file is absent, no entry is
created and the row's \code{.get(..., \{\})} returns a bare \code{\{\}}. 
Both end up empty, but only the domain case leaves a record and an empty \code{action_names} set means \emph{every} emitted action looks illegal,
silently depressing that domain's grounding numbers. The only trace is a warning, which is exactly what input~3 is for.
\end{bmcavbox}

\paragraph{\texorpdfstring{\texttt{warnings\_out}}{warnings\_out}: the diagnostic accumulator}

\code{warnings_out} is a single \code{list[dict]} created once at the top of the
analysis run and threaded by reference into every stage that can encounter a
malformed or missing input in \code{index_pddl_context} for both PDDL parsers and
\code{compute_row_metrics}. Each stage appends through one helper:

\begin{lstlisting}
def add_warning(warnings_out, warning_type, message, **extra):
    warnings_out.append({"type": warning_type, "message": message, **extra})
    # **extra carries locating context: run_id, domain, difficulty, problem, path
\end{lstlisting}

Its purpose is to replace scattered \code{print}/\code{logging} calls with one
structured, serialisable trace that travels with the results: the full list is
embedded in \code{report.json} under the \code{warnings} key, so a notebook or CI
job can filter by \code{type} and locate the offending artifact without
re-running anything. Typical entries are \code{missing_tasks_dir},\code{missing_domain_pddl}, \code{missing_pddl_context},
\code{domain_parser_error}, \code{problem_parser_error}. Rows with incomplete context are \emph{not} dropped they still get NaN metric values so the
instance count stays honest while the warning records why those NaNs exist.

\begin{bmcavbox}{One diagnostic path bypasses \code{warnings\_out}}
The simulator is the exception. If \code{compute_precondition_metrics} throws, it
falls back to a NaN result and reports through Python's stdlib
\code{warnings.warn("[simulator] ...")} \emph{not} through
\code{add_warning}. So a simulator crash does not appear in the structured
\code{warnings} list in the JSON report it surfaces only on stderr. Worth
unifying if the report is to be the single audit log.
\end{bmcavbox}

% ---------------------------------------------------------------------------
\subsubsection{The merge: three inputs into one widened frame}
% ---------------------------------------------------------------------------

\code{compute_row_metrics} iterates \code{df_raw} once. For each row it resolves
\code{d_info}/\code{p_info} (input~2), runs the three metric families against the
row's \code{_actions} and attempt lists, emits any incompleteness into
\code{warnings_out} (input~3). The per-family results are collected and
concatenated columnwise back onto the copied skeleton.

\begin{lstlisting}
df = df_raw.copy()
for _, row in df.iterrows():
    d_info  = domain_info.get(row["Domain"], {})
    p_info  = problem_info.get((row["Domain"], row["Difficulty"], row["Problem"]), {})
    if not (d_info.get("action_names") and p_info.get("objects")):
        add_warning(warnings_out, "missing_pddl_context", ...)        # input 3
    halluc_rows.append(      compute_hallucination_metrics(actions, d_info, p_info))
    precondition_rows.append(compute_precondition_metrics(actions, d_info, p_info))  # PDDLSimulator
    cot_rows.append(         ... compute_cot_alignment_for_attempt(...) ...)

df_metrics = pd.concat([df, halluc_df, precondition_df, cot_df], axis=1)
df_metrics["_iter1"]        = (df.Valid == True) & (df.Iterations == 1)   # -> FASR
df_metrics["_iwsr_contrib"] = 1/Iterations if Valid else 0                # -> IWSR
\end{lstlisting}

The result \code{df_metrics} is the single source of truth every
\code{groupby} in \code{build_aggregate_tables} and every plot reads from. Note
the grain: the hallucination and executability columns describe the
\emph{representative} (last non-empty) plan, while the CoTaL and efficiency
signals reach back into the preserved attempt lists.

% ---------------------------------------------------------------------------
\subsubsection{Data classes that make the metrics possible}
% ---------------------------------------------------------------------------

Two classes carry state that a flat function could not  expanded in full in Section~\ref{sec:sim}.

\paragraph{PDDLSimulator (in the analysis script).} A forward STRIPS +
simple-numeric simulator written in pure Python so the analysis layer needs no
VAL subprocess. Instantiated \emph{once per (problem, plan)} inside
\code{compute_precondition_metrics}. Its decisive feature is \code{prop_hist}: a
list of \code{frozenset} snapshots of the propositional state after every
applied step. That history is what lets the framework at the first failing
action, ask \emph{``was this precondition ever true earlier?''} that is the question
separating a sequencing error (true before, deleted too early) from a state
fabrication (never true). Without the per-step history there is no PAS and no
temporal distance.

\paragraph{RunMetrics (\texttt{evaluators/metrics.py}, \texttt{@dataclass}).} The
validator-side metric bundle written at scoring time: \code{validity_at_1},
\code{validity_at_k}, \code{repair_success}, \code{iterations_to_valid},
\code{plan_length}, \code{error_type}, \code{hit_iteration_limit}. The analysis
script re-derives most of these from the scored JSON but this dataclass is the
contract guaranteeing every model is scored on the same basis, independent of
backend.

The error vocabulary is fixed in \code{evaluators/error_taxonomy.py} as an
\code{ErrorType} enum, partitioned into \code{LOGICAL_ERROR_TYPES} (empty plan,
syntax, unknown action, invalid precondition, unsatisfied goal) versus
\code{TECHNICAL_ERROR_TYPES} (parse error, timeout, validator crash/unavailable).
That split keeps infrastructure failures from being charged against a model's
planning score.

% ---------------------------------------------------------------------------
\subsubsection{\texorpdfstring{\texttt{PDDLSimulator}}{PDDLSimulator} in depth}
\label{sec:sim}
% ---------------------------------------------------------------------------

Almost every execution-side discriminator: executability ratio, the
sequencing-vs-fabrication split, PAS, and mean temporal distance is produced
by one small class. It is a \textbf{forward STRIPS + simple-numeric simulator
written in pure Python}, deliberately so: the analysis layer can re-derive these
signals without spawning a VAL subprocess or shipping a compiled binary. It
reproduces just enough of the validator's logic for the propositional and
numeric-inequality fragments the benchmark uses by trading completeness for
portability.

\paragraph{What it is for.} The class does three things: it \textbf{holds a
world state} (a set of true propositions plus a dict of numeric fluents), it
\textbf{walks the plan one action at a time} checking each action's preconditions before applying its effects it
\textbf{saves a frozen snapshot of the state after every applied step} which is the decisive feature. That
snapshot history is what makes the precondition-awareness metrics computable. It
is the basis for four metrics: executability ratio (how far the plan runs before
the first break), the sequencing-error / state-fabrication counts, PAS (their
ratio), and mean temporal distance. It targets \emph{forward},
\emph{propositional} and \emph{simple-numeric} domains , STRIPS add/delete
effects plus \code{increase}/\code{decrease} fluents and
$\geq\ \leq\ >\ <$ comparisons and is invoked exactly once per
\code{(problem, plan)} pair.

\begin{figure}[h]
\centering
\resizebox{\linewidth}{!}{%
\begin{tikzpicture}[
  font=\ttfamily\scriptsize,
  snap/.style={draw=bmblue,fill=white,rounded corners=3pt,minimum width=11mm,
               minimum height=8mm,line width=1pt,text=bmblue,font=\ttfamily\bfseries\small},
  fail/.style={draw=bmred,fill=bmred!8,rounded corners=3pt,minimum width=12mm,
               minimum height=8mm,line width=1pt,text=bmred,font=\ttfamily\bfseries\small},
  fwd/.style={-{Stealth[length=5pt]},bmblue,line width=1pt},
  back/.style={-{Stealth[length=5pt]},bmred,line width=1pt,dashed},
  node distance=11mm,
]
\node[snap] (s0) {S0};
\node[snap,right=of s0] (s1) {S1};
\node[snap,right=of s1] (s2) {S2};
\node[snap,right=of s2] (s3) {S3};
\node[fail,right=18mm of s3] (fk) {$a_k$ $\times$};

\draw[fwd] (s0) -- node[above=-1pt,text=muted,font=\ttfamily\tiny]{$a_0$} (s1);
\draw[fwd] (s1) -- node[above=-1pt,text=muted,font=\ttfamily\tiny]{$a_1$} (s2);
\draw[fwd] (s2) -- node[above=-1pt,text=muted,font=\ttfamily\tiny]{$a_2$} (s3);
\draw[fwd] (s3) -- node[above=-1pt,text=muted,font=\ttfamily\tiny]{$\dots a_k$} (fk);

\draw[back] (fk.north) to[out=120,in=60] node[above,text=bmred,font=\ttfamily\tiny]
  {temporal distance $= k-0$} (s0.north);

\node[text=muted,font=\sffamily\scriptsize,anchor=west] at ([yshift=-9mm]s0.south west)
  {prop\_hist $= [\,$frozenset(S0), frozenset(S1), $\dots\,]$ --- one snapshot per applied step};
\end{tikzpicture}%
}
\caption{State evolution. At the first failing action $a_k$, each unmet
precondition atom is traced back through \texttt{prop_hist}: found earlier
$\rightarrow$ \textcolor{bmblue}{sequencing error} (record $k-\text{last}$ as
temporal distance); never found $\rightarrow$ \textcolor{bmred}{state
fabrication}.}
\label{fig:sim}
\end{figure}

\paragraph{State, snapshots, and the step loop.} The constructor seeds the world
from the parsed problem and takes the first snapshot before any action runs;
\code{prop_hist} then grows by one frozen set per applied step.

\begin{lstlisting}
def __init__(self, d_info, p_info):
    self.prop      = set(p_info.get("init_atoms", set()))   # true propositions
    self.num       = dict(p_info.get("init_numeric", {}))   # numeric fluents
    self.prop_hist = [frozenset(self.prop)]                 # snapshot S0 = initial state
    self.step      = 0
\end{lstlisting}

For each action, \code{compute_precondition_metrics} grounds the schema, checks
preconditions and only if they hold applies the effects and appends a
new snapshot. The walk \textbf{stops at the first failing action} (a
\code{break}), which makes ``how far did it get'' a meaningful prefix length.

\begin{lstlisting}
for step, action_str in enumerate(actions):
    name, args = parse_action(action_str)
    schema = schemas.get(name)
    if name is None or schema is None:        # unparseable / unknown action: skipped
        continue
    if len(schema["params"]) != len(args):    # wrong arity -> fabrication, stop
        first_failure = step; fab_errors += 1; break
    ok, failed_prop, failed_num = sim.check_preconditions(schema["prec_raw"], params, args)
    if not ok:
        first_failure = step
        for atom in failed_prop:
            last = sim.last_step_atom_true(atom)         # search the snapshot history
            if last >= 0: seq_errors += 1; 
            temporal_distances.append(step - last)
            else:         fab_errors += 1
        fab_errors += len(failed_num); break
    sim.apply_effects(schema["eff_raw"], params, args)   # mutate state + append snapshot
\end{lstlisting}

\paragraph{The classification logic, atom by atom.} At the failing action, each
unsatisfied propositional precondition is the unit of analysis.
\code{last_step_atom_true(atom)} scans \code{prop_hist} newest-to-oldest and
returns the latest snapshot index at which the atom held or \code{-1} if it never did.

\begin{lstlisting}
def last_step_atom_true(self, atom):
    for index in range(len(self.prop_hist) - 1, -1, -1):
        if atom in self.prop_hist[index]:
            return index
    return -1
\end{lstlisting}

That single lookup is the whole basis of the sequencing/fabrication distinction.
If the atom was true at some earlier snapshot, the model had the prerequisite
and then ordered actions so as to destroy it we have a \textbf{sequencing error},
and the gap $\text{step}-\text{last}$ is recorded as a temporal distance. If the
atom was never true, the model required a world state that never existed we have a
\textbf{state fabrication}. Effects are applied in STRIPS order (deletes before
adds), so the snapshots faithfully reflect what each action left behind. The
metrics fall straight out of these counts:

\begin{bmformulabox}
\ttfamily\small
executability_ratio = first_failure / plan_len\\[2pt]
PAS = sequencing_errors / (sequencing_errors + state_fabrications)\\[2pt]
mean_temporal_distance = mean(step - last_true_step  for each sequencing error)
\end{bmformulabox}

\begin{bmcavbox}{The counts describe the first failure, not the whole plan}
Because the loop \code{break}s at the first failing action, the
sequencing/fabrication counts and the temporal distances come from that
\emph{one} action's unmet preconditions and PAS is a property of where the plan
first breaks not a tally over the entire sequence. This is intended (it
localises the failure) but it means a plan can accrue at most one failure event
in these columns.
\end{bmcavbox}

\begin{bmcavbox}{Unknown actions are skipped, which inflates executability}
An action whose name is not in the schema is \code{continue}d over it neither
advances the state nor counts as a failure here (name-level hallucination is
scored separately). So a plan made entirely of hallucinated action names never
triggers a precondition failure and records \code{executability_ratio = 1.0}.
Always read ER alongside the hallucination columns. Relatedly, because skipped
actions advance the plan index but not the snapshot index, temporal distance is
measured in generated-plan positions relative to the snapshot where the atom
last held the two clocks coincide only when every prior action was a known,
applied schema.
\end{bmcavbox}

% ---------------------------------------------------------------------------
\subsubsection{From data to metrics to pictures --- why each stage exists}
% ---------------------------------------------------------------------------

The architecture separates concerns on purpose, so each is independently
testable.

\begin{table}[h]\small
\begin{tabularx}{\linewidth}{@{}l l >{\raggedright\arraybackslash}X@{}}
\toprule
\textbf{Stage} & \textbf{Function} & \textbf{Purpose / why it is separate}\\
\midrule
Load & \code{load_artifact_rows} &
Decouples the filesystem walk from metric logic; everything below receives a
uniform table.\\[3pt]
Row metrics & \code{compute_row_metrics} &
One pass attaches every per-instance metric; private \code{_actions} /
\code{_cot_text} stay alive only here.\\[3pt]
Aggregate & \code{build_aggregate_tables} &
All \code{groupby} logic in one place so the JSON report can embed tables whether
or not plots are drawn.\\[3pt]
Visualise & \code{maybe_make_plots} &
Reads only the aggregate tables; registering each figure embeds its path +
description into the report.\\[3pt]
Assemble & \code{build_model_payloads} &
Re-keys everything by \emph{model} so a consumer can pull one model in O(1), and
attaches the capability-profile classification.\\
\bottomrule
\end{tabularx}
\end{table}

\begin{bmcavbox}{A real artefact to know about}
Because \code{compute_row_metrics} keeps the last non-empty plan as the
representative for the instance-level hallucination/executability columns, those
columns describe the \emph{final} attempt. Metrics that must see all attempts
(FASR, IWSR, CoTaL per-iteration, the iteration profile) read the preserved
\code{_*_attempts} lists instead. When reading a plot, always check which grain
it is on --- this is spelled out per metric in the following sections.
\end{bmcavbox}

